\chapter{2006年全国I卷（理）}

\section{单选题}

\begin{problem}
设集合 \(M = \{x \mid x^2 - x < 0\}\)，\(N = \{x \mid |x| < 2\}\)，则（\quad）

\choices
    {\(M \cap N = \varnothing\)}
    {\(M \cap N = M\)}
    {\(M \cup N = M\)}
    {\(M \cup N = \R\)}
\end{problem}

\begin{problem}
已知函数 \(y = \e^x\) 的图象与函数 \(y = f(x)\) 的图象关于直线 \(y = x\) 对称，则（\quad）

\choices
    {\(f(2x) = \e^{2x}\)（\(x \in \R\)）}
    {\(f(2x) = \ln 2\cdot\ln x\)（\(x > 0\)）}
    {\(f(2x) = 2\e^x\)（\(x \in \R\)）}
    {\(f(2x) = \ln x + \ln 2\)（\(x > 0\)）}
\end{problem}

\begin{problem}
双曲线 \(mx^2 + y^2 = 1\) 的虚轴长是实轴长的 \(2\) 倍，则 \(m=\)（\quad）

\choices
    {\(-\frac{1}{4}\)}
    {\(-4\)}
    {\(4\)}
    {\(\frac{1}{4}\)}
\end{problem}

\begin{problem}
如果复数 \((m^2 + \i)(1 + m\i)\) 是实数，则实数 \(m=\)（\quad）

\choices
    {\(1\)}
    {\(-1\)}
    {\(\sqrt{2}\)}
    {\(-\sqrt{2}\)}
\end{problem}

\begin{problem}
函数 \(f(x)=\tan\left(x+\frac{\pi}{4}\right)\) 的单调增区间为（\quad）

\choices
    {\(\left(k\pi-\frac{\pi}{2},k\pi+\frac{\pi}{2}\right)\)（\(k\in\Z\)）}
    {\((k\pi,(k+1)\pi)\)（\(k\in\Z\)）}
    {\(\left(k\pi-\frac{3\pi}{4},k\pi+\frac{\pi}{4}\right)\)（\(k\in\Z\)）}
    {\(\left(k\pi-\frac{\pi}{4},k\pi+\frac{3\pi}{4}\right)\)（\(k\in\Z\)）}
\end{problem}

\begin{problem}
\(\triangle ABC\) 的内角 \(A\)，\(B\)，\(C\) 的对边分别为 \(a\)，\(b\)，\(c\). 若 \(a\)，\(b\)，\(c\) 成等比数列，且 \(c=2a\)，则 \(\cos B=\)（\quad）

\choices
    {\(\frac14\)}
    {\(\frac34\)}
    {\(\frac{\sqrt2}{4}\)}
    {\(\frac{\sqrt2}{3}\)}
\end{problem}

\begin{problem}
已知各顶点都在一个球面上的正四棱柱高为 \(4\)，体积为 \(16\)，则这个球的表面积是（\quad）

\choices
    {\(16\pi\)}
    {\(20\pi\)}
    {\(24\pi\)}
    {\(32\pi\)}
\end{problem}

\begin{problem}
抛物线 \(y=-x^2\) 上的点到直线 \(4x+3y-8=0\) 距离的最小值是（\quad）

\choices
    {\(\frac43\)}
    {\(\frac75\)}
    {\(\frac85\)}
    {\(3\)}
\end{problem}

\begin{problem}
设平面向量 \(\bs{a}_1\)，\(\bs{a}_2\)，\(\bs{a}_3\) 的和为 \(\bs{a}_1+\bs{a}_2+\bs{a}_3=0\). 如果向量 \(\bs{b}_1\)，\(\bs{b}_2\)，\(\bs{b}_3\) 满足 \(|\bs{b}_i|=2|\bs{a}_i|\)，且 \(\bs{a}_i\) 顺时针旋转 \(30^{\circ}\) 后与 \(\bs{b}_i\) 同向，其中 \(i=1\)，\(2\)，\(3\)，则（\quad）

\choices
    {\(-\bs{b}_1+\bs{b}_2+\bs{b}_3=0\)}
    {\(\bs{b}_1-\bs{b}_2+\bs{b}_3=0\)}
    {\(\bs{b}_1+\bs{b}_2-\bs{b}_3=0\)}
    {\(\bs{b}_1+\bs{b}_2+\bs{b}_3=0\)}
\end{problem}

\begin{problem}
设 \(a_n\) 是公差为正数的等差数列，若 \(a_1+a_2+a_3=15\)，\(a_1a_2a_3=80\)，则 \(a_{11}+a_{12}+a_{13}=\)（\quad）

\choices
    {\(120\)}
    {\(105\)}
    {\(90\)}
    {\(75\)}
\end{problem}

\begin{problem}
用长度分别为 \(2\)，\(3\)，\(4\)，\(5\)，\(6\)（单位：\(\mathrm{cm}\)）的 \(5\) 根细木棒围成一个三角形（允许连接，但不允许折断），能够得到的三角形的最大面积为（\quad）

\choices
    {\(8\sqrt5\mathrm{~cm}^{2}\)}
    {\(6\sqrt{10}\mathrm{~cm}^{2}\)}
    {\(3\sqrt{55}\mathrm{~cm}^{2}\)}
    {\(20\mathrm{~cm}^{2}\)}
\end{problem}

\begin{problem}
设集合 \(I=\{1,2,3,4,5\}\). 选择 \(I\) 的两个非空子集 \(A\) 和 \(B\)，要使 \(B\) 中最小的数大于 \(A\) 中最大的数，则不同的选择方法共有（\quad）

\choices
    {\(50\) 种}
    {\(49\) 种}
    {\(48\) 种}
    {\(47\) 种}
\end{problem}

\section{填空题}

\begin{problem}
已知正四棱锥的体积为 \(12\)，底面对角线长为 \(2\sqrt6\)，则侧面与底面所成的二面角等于\fillinblank{}.
\end{problem}

\begin{problem}
设 \(z=2y-x\)，式中变量 \(x\)，\(y\) 满足下列条件：\(\begin{cases}
2x-y\geqslant-1,\\
3x+2y\leqslant23,\\
y\geqslant1.
\end{cases}\) 则 \(z\) 的最大值为\fillinblank{}.
\end{problem}

\begin{problem}
安排 \(7\) 位工作人员在 \(5\) 月 1 日至 \(5\) 月 7 日值班，每人值班一天，其中甲、乙二人都不安排在 5 月 1 日和 2 日，不同的安排方法共有\fillinblank{}种.（用数字作答）
\end{problem}

\begin{problem}
函数 \(f(x)=\cos(\sqrt3x+\varphi)\)（\(0<\varphi<\pi\)）. 若 \(f(x)+f'(x)\) 是奇函数，则 \(\varphi=\)\fillinblank{}.
\end{problem}

\section{解答题}

\begin{problem}
\(\triangle ABC\) 的三个内角为 \(A\)，\(B\)，\(C\)，求当 \(A\) 为何值时，\(\cos A+2\cos\frac{B+C}{2}\) 取得最大值，并求出这个最大值.
\end{problem}

\begin{problem}
\(A\)，\(B\) 是治疗同一种疾病的两种药，用若干试验组进行对比试验. 每个试验组由 \(4\) 只小白鼠组成，其中 \(2\) 只服用 \(A\)，另 \(2\) 只服用 \(B\)，然后观察疗效. 若在一个试验组中，服用 \(A\) 有效的小白鼠的只数比服用 \(B\) 有效的多，就称该试验组为甲类组. 设每只小白鼠服用 \(A\) 有效的概率为 \(\frac23\)，服用 \(B\) 有效的概率为 \(\frac12\).

\begin{enumerate}
\item 求一个试验组为甲类组的概率.
\item 观察 \(3\) 个试验组，用 \(\xi\) 表示这 \(3\) 个试验组中甲类组的个数，求 \(\xi\) 的分布列和数学期望.
\end{enumerate}
\end{problem}

\begin{problem}
如图，\(l_1\)，\(l_2\) 是互相垂直的异面直线，\(MN\) 是它们的公垂线段. 点 \(A\)，\(B\) 在 \(l_1\) 上，\(C\) 在 \(l_2\) 上，\(AM=MB=MN\).

\begin{enumerate}
\item 证明：\(AC\perp NB\).
\item 若 \(\angle ACB=60^{\circ}\)，求 \(NB\) 与平面 \(ABC\) 所成角的余弦值.
\end{enumerate}

\bitmapfigure[float=inline,width=5.4cm]{img/2006/national_paper_1_science/q19_fig1.png}
\end{problem}

\begin{problem}
在平面直角坐标系 \(xOy\) 中，有一个以 \(F_1(0,-\sqrt3)\) 和 \(F_2(0,\sqrt3)\) 为焦点、离心率为 \(\frac{\sqrt3}{2}\) 的椭圆，设椭圆在第一象限的部分为曲线 \(C\)，动点 \(P\) 在 \(C\) 上，\(C\) 在点 \(P\) 处的切线与 \(x\)，\(y\) 轴的交点分别为 \(A\)，\(B\)，且向量 \(\overrightarrow{OM}=\overrightarrow{OA}+\overrightarrow{OB}\). 求：

\begin{enumerate}
\item 点 \(M\) 的轨迹方程.
\item \(|\overrightarrow{OM}|\) 的最小值.
\end{enumerate}
\end{problem}

\begin{problem}
已知函数 \(f(x)=\frac{1+x}{1-x}\e^{-ax}\).

\begin{enumerate}
\item 设 \(a>0\)，讨论 \(y=f(x)\) 的单调性.
\item 若对任意 \(x\in(0,1)\) 恒有 \(f(x)>1\)，求 \(a\) 的取值范围.
\end{enumerate}
\end{problem}

\begin{problem}
设数列 \(a_n\) 的前 \(n\) 项的和 \(S_n=\frac43a_n-\frac13\times2^{n+1}+\frac23\)，\(n=1\)，\(2\)，\(3\)，\(\ldots\).

\begin{enumerate}
\item 求首项 \(a_1\) 与通项 \(a_n\).
\item 设 \(T_n=\frac{2^n}{S_n}\)，\(n=1\)，\(2\)，\(3\)，\(\ldots\)，证明：\(\sum_{i=1}^nT_i<\frac32\).
\end{enumerate}
\end{problem}
